% Part: incompleteness
% Chapter: theories-computability
% Section: intepretability

\documentclass[../../../include/open-logic-section]{subfiles}

\begin{document}

\olfileid{inc}{tcp}{itp}

\olsection{هغه تيورۍ ناپرېکړه‌وړې دي چې~$\Th{Q}$ پکې تفسيرېدونکې وي}

دا پايلې لا نورې هم پياوړې کولے شو۔ په غيرصوري ډول، په بله
ژبه~$\Lang{L_2}$ کښې د يوې ژبې~$\Lang{L_1}$ تفسيرول دا غواړي چې د
$\Lang{L_1}$ کائنات، اړيکيزې نښې او تابعي نښې په~$\Lang{L_2}$ کښې
د !!{formula}s په وسيله تعريف شي۔ که څه هم دلته ورته وخت نۀ ورکوو،
دا تعريف په کره بڼه جوړېدے شي۔

\begin{thm}
  فرض کړئ~$\Th{T}$ په داسې ژبه کښې يوه تيوري ده چې د حساب ژبه پکې
  داسې تفسيرېدے شي چې~$\Th{T}$ د~$\Th{Q}$ له تفسير سره سازګاره وي۔
  بيا~$\Th{T}$ ناپرېکړه‌وړې ده۔ که~$\Th{T}$ د~$\Th{Q}$ د بديهي
  اصولو تفسير ثابتوي، نو د~$\Th{T}$ هېڅ سازګاره غځونه د پرېکړې وړ
  نۀ ده۔
\end{thm}

ثبوت د تېرې قضيې د ثبوت يوازې يو وړوکے بدلون دے؛ د يوې خلاف بېلګې
په کارولو د~$\Th{Q}$ او~$\Th{\bar Q}$ بېلوونکې ترلاسه کېدے شي۔
$\Th{ZFC}$، يعنې د انتخاب له بديهي اصل سره د زر مېلو--فرېنکل سټ
تيوري، داسې بديهي بنسټ ګڼلے شو چې د عادي رياضياتو ډېره برخه پکې
ترسره کېدے شي۔ په~$\Th{ZFC}$ کښې طبيعي عددونه تعريفولے شو، او د
دې تفسير له مخې د~$\Th{Q}$ بديهي اصول رښتيا دي۔ نو لرو:

\begin{cor}
د~$\Th{ZFC}$ هېڅ د پرېکړې وړ غځونه نشته۔
\end{cor}

\begin{cor}
د~$\Th{ZFC}$ هېڅ بشپړه، سازګاره او په محاسبوي ډول
!!{axiomatizable} غځونه نشته۔
\end{cor}

د~$\Th{ZFC}$ ژبه يوازې يوه دوه‌ځاييزه اړيکه~$\in$ لري۔ (په حقيقت
کښې ان عينيت ته هم اړتيا نشته۔) نو لرو:

\begin{cor}
د هرې هغې ژبې دپاره لومړۍ درجې منطق ناپرېکړه‌وړ دے چې يوه
دوه‌ځاييزه اړيکيزه نښه ولري۔
\end{cor}

دا پايله هرې هغې ژبې ته هم غځېږي چې دوه يوځاييزې تابعي نښې ولري،
ځکه د هغو په وسيله يوه دوه‌ځاييزه اړيکيزه نښه شبيه کېدے شي۔ همدا
اوس يادې شوې پايلې کره پولې دي: په حقيقت کښې د هرې هغې ژبې دپاره
لومړۍ درجې منطق د پرېکړې وړ دے چې يوازې \emph{يوځاييزې} اړيکيزې
نښې او تر ډېره يوه \emph{يوځاييزه} تابعي نښه ولري۔

يو بل په زړه پورې حقيقت۔ پوهېږو چې په معياري مدل کښې د ژبې
$\Obj 0$، $'$، $+$، $\times$، $<$ د رښتينو جملو سټ ناپرېکړه‌وړ
دے۔ په حقيقت کښې~$<$ د نورو نښو په وسيله تعريفولے شو، او بيا~$+$
د~$\times$ او~$'$ په وسيله تعريفولے شو۔ نو د ژبې~$\Obj 0$، $'$،
$\times$ د رښتينو جملو سټ ناپرېکړه‌وړ دے۔ له بلې خوا، پرېسبورګر
ښودلې ده چې د ژبې~$\Obj 0$، $'$، $+$ د هغو جملو سټ چې د حساب په
معياري تفسير کښې رښتيا دي، د پرېکړې وړ دے۔
\textbf{د ژباړې د نسخې سمون (OLCMP-096):} سرچينې وروستے صدق «د
حساب په ژبه کښې» بللے و؛ ژبه يوازې نحو ټاکي، نو دلته هغه معياري
تفسير په څرګنده ياد شوے چې د پرېسبورګر د پرېکړې د پايلې موخه ده۔
خو دا کړنلاره له محاسبوي پلوه عملي نۀ ده۔

\end{document}
